\documentclass[11pt]{article}
\usepackage[margin=1in]{geometry}
\usepackage{booktabs}
\usepackage{array}
\usepackage{amsmath}
\usepackage{url}

\title{\vspace{-2.5em}Row and Permit Attrition Across the DMR Filter Pipeline\\
FY2009 vs.\ FY2025}
\author{}
\date{\today}

\begin{document}
\maketitle
\vspace{-2.5em}

\section*{Scope and data}
This brief tracks how many \textbf{observations} (DMR rows) and how many
\textbf{distinct permits} survive each step of the DMR filter pipeline in
\texttt{dmr analysis/}, comparing FY2009 against FY2025. The starting point
(\textbf{Stage 0}) is the \emph{raw, completely unfiltered} DMR file exactly
as delivered by EPA ECHO --- every permit (major and minor, individual and
general), every parameter, feature type, and monitoring location, with no
restriction of any kind. Four filtering stages are then applied in sequence,
narrowing that raw universe down to a single pollutant/monitoring-location/
limit-type slice used for TSS effluent compliance analysis:

\begin{enumerate}
  \item \textbf{Major-Individual} --- restrict to permits that are
        individual (\texttt{PERMIT\_TYPE\_CODE = "NPD"}) and were flagged
        major in at least one permit version. No parameter, feature-type, or
        monitoring-location restriction yet.
  \item \textbf{+TSS (00530)} --- restrict to \texttt{PARAMETER\_CODE =
        '00530'} (Total Suspended Solids).
  \item \textbf{+Effluent Gross} --- restrict to
        \texttt{MONITORING\_LOCATION\_CODE IN ('1','EG')} (both map to
        "Effluent Gross" per EPA's \texttt{REF\_MONITORING\_LOCATION} table).
  \item \textbf{+C1/Q1} --- restrict to \texttt{LIMIT\_VALUE\_TYPE\_CODE IN
        ('C1','Q1')} (the concentration- and mass-average value slots).
\end{enumerate}

Each stage is built by \texttt{dmr analysis/filter\_dmr\_*.R} (parameterized
by fiscal year) and self-verifies column counts and category cardinality
before writing its output; see the scripts for the exact filter logic. Stage
0's summary is computed directly via DuckDB SQL aggregation rather than
loaded into R (the raw file is 11--27M rows, too large to safely read in full
on this machine); see \texttt{scripts/summary/build\_dmr\_raw\_summary.R}.
All figures below are measured directly from the raw file and the four
filtered CSVs for each fiscal year, not estimated.

\section*{Underlying numbers}

\begin{table}[h!]
\centering
\begin{tabular}{lrrrr}
\toprule
 & \multicolumn{2}{c}{FY2009} & \multicolumn{2}{c}{FY2025} \\
\cmidrule(lr){2-3} \cmidrule(lr){4-5}
Stage & Rows & Permits & Rows & Permits \\
\midrule
0. Raw / All Permits  & 11{,}077{,}254 & 44{,}230 & 26{,}869{,}750 & 89{,}954 \\
1. Major-Individual   & 4{,}015{,}793 & 6{,}555 & 4{,}703{,}897 & 6{,}701 \\
2. +TSS (00530)        &   440{,}369 & 6{,}426 &   489{,}033 & 6{,}572 \\
3. +Effluent Gross     &   328{,}296 & 6{,}091 &   336{,}437 & 6{,}117 \\
4. +C1/Q1              &    80{,}196 & 4{,}974 &    75{,}818 & 4{,}880 \\
\bottomrule
\end{tabular}
\caption{Row and permit counts at each filter stage, FY2009 vs.\ FY2025.}
\label{tab:funnel}
\end{table}

\end{document}
